\chapter{Literature Study}
\label{chapter:literature-study}
The following chapter presents a literature study that goes through both related work and background necessary for understanding the context of this thesis. The literature study has been conducted as a comprehensive literature review \cite{williams2018}. It starts simple with defining cyber grooming, then, describing different detection methods for prevention, and how evidence in a cyber grooming case is located. Furthermore, it describes \gls{llm} and explores the way \gls{llm}s and \gls{ai} can be utilised in cyber grooming evidence localisation. Lastly, the chapter closes with the research gap that this thesis aims to cover. 

\section{Cyber grooming is a process}
\label{literature-study:cyber-grooming}
In 2002, Gillespie \cite{gillespie2002} defined grooming in general, as the process in which an offender works on building a relationship with a child to gain their trust and confidence, making the child more compliant to be exploited. Some later definitions of grooming vary in detail, but grooming got a new category when internet became one of the most common platforms to start the process: cyber grooming. Cyber grooming definitions also vary, but the process with an intention to exploit minors online is prevalent \cite{klicksafe2026, nspcc2026, childsafenet2026, canada2026}. The Cyber Security Authority \cite{CSA2026} adds that the offender may even pretend to be child in order to gain their victim's trust and confidence. What gives cyber grooming its own category is the platform in which it is conducted. An offender operates in a setting with low risk, where they can engage with a child without letting the child know who they are actually talking to. The offender can easily disengage and disappear without being identified, and for law enforcement this is difficult to track. In a systematic review, Schittenhelm et al. \cite{schittenhelm2025} found that reported prevalence of cyber grooming varies between studies, but that psychological harm to victims is serious, consistent, and evident. 

Cyber grooming is a process that unfolds over time, not an act contained in a single chat message. O'Connell \cite{oconnell2003} proposed a model in 2003 that described six stages of cyber grooming: friendship forming, relationship forming, risk assessment, exclusivity, a sexual stage, and a conclusion in which the offender works towards an offline meeting. O'Connell noted that the stages do not need to be in order and that their extents vary between conversations.
Later empirical work has studied and questioned how closely real conversations follow O'Connell's model. Black et al. \cite{black2015} analysed 44 transcripts from convicted online offenders to assess the stages in the model. They found evidence of the model's strategies, in different order and with different timings. Lorenzo-Dus et al. \cite{lorenzodus2016} examined the perverted justice \cite{pervertedjustice} corpus to research the characteristics of online grooming. They proposed a new model of online grooming discourse: deceptive trust development, sexual gratification, compliance testing, and isolation. On this account, the offenders grooming strategies occur in parallel.

In research "grooming" is narrowed to the slow process of building trust before anything sexual happens. However, not all offenders follow the gradual process the above mentioned models describe. Kloess et al. \cite{kloess2019} analysed five offenders with 29 transcripts in police cases involving real victims. Two of the offenders took an indirect approach like the models described above. The other three took a more direct approach, moving to the sexual content quickly and skipping the trust-building stages that define grooming. In a recent systematic review, Moosburner et al. \cite{moosburner2026} likewise distinguish between the direct, sexually forward strategies, from the slower, relationship-building ones. 

For evidence localisations, these findings have a practical consequence. If the grooming behaviour is spread across a conversation, amid ordinary talk and varies in style between offenders, then the evidentiary value of a conversation is not evenly distributed. Some segments of a conversation may carry much more grooming content than others, and which ones depend on the offender's strategy. A verdict for the whole conversation therefore says little about where an investigator should look. This motivates treating evidence localisation as a ranking of conversational segments, rather than classifying conversations. It suggests evidence should be recognised in several kinds such as trust building, sexual content, isolation, and risk assessment. 

\section{Datasets, sources and their limitations}
\label{literature-study:cyber-grooming-datasets}
Most research on grooming detection and cyber grooming in general depends on a small number of public datasets. Real case material is protected by law and are rarely released. This thesis uses two publicly accessed datasets, the \gls{pj} archive and the PAN12 dataset, which are described in further detail in \autoref{chapter:datasets}. This section does not describe the datasets themselves, but reviews what earlier studies have found about them and what those findings mean for research built on them.

Some studies have examined how the \gls{pj} and PAN12 datasets were built. Gupta et al. \cite{gupta2012} pointed out that the \gls{pj} logs were gathered from different sources, contain the volunteers' comments mixed into the chat text, and that the logs cannot be checked for completeness. PAN12 was designed for a shared task \cite{inches2012}. One of the tasks asked participants to identify the lines that showed most typical grooming behaviour. Vogt et al. \cite{vogt2021} surveyed the available datasets, documented their limitations for early detection, and showed that choices made in constructing a dataset affect the results reported on it. They also noted that PAN12 reuses \gls{pj} material, an overlap that this thesis must handle for its cross-dataset evaluation (see \autoref{chapter:datasets}).

One question raised in the literature is whether these datasets represents grooming at all, since both datasets contain conversations with decoys rather than children. Ringenberg \cite{ringenberg2021} points out that decoys and police officers cannot initiate contact, suggest a meeting, or start a sexual conversation. However, because their goal is to gather evidence, they often respond positively where a real child might hesitate. 

Studies based on real transcripts between offenders and victims do exist, but their data is not public \cite{kloess2017, kloess2019}. 

A response to the data scarcity of cyber grooming conversations has been to generate the data. Zhang et al. \cite{xhang2024} and An et al. \cite{an2026} used \gls{llm}s to generate grooming conversations. %write more about how they did this, and if their results show promise or not. 

\section{Automated detection methods}
Automated grooming detection is an active research area that has moved from rules and hand-crafted features, through classical and deep machine learning, to \gls{llm}s. 

Early systems drew directly on grooming theory. McGhee et al. \cite{mcghee2011} built ChatCoder, which used phrase matching and rules to label individual chat lines with categories derived from communication theory. Machine learning classifiers trained on the same task did not improve on the rule-based system's accuracy of 68 percent. The PAN-2012 competition then established the standard setting for the field \cite{inches2012}. The winning system used a two-step approach: it first filtered out conversations unlikely to involve a predator, then identified the predator within the remaining ones \cite{villatoro2012}. Later systems followed the same two-stage design. Feature-based work continued after the competition. Parapar et al. \cite{parapar2014} combined psycholinguistic, content-based and chat-behaviour features with \gls{svm}. Borj and Bours \cite{borjbours2019} proposed predatory conversation detection using linguistic and behavioural features, and Bours and Kulsrud \cite{bours2019} applied \gls{svm}s with lexical and stylistic features to separate grooming conversations from ordinary ones.

Later work improved classification with ensembles, embeddings and neural networks \cite{fauzi2020, borj2021, borj2023}. Deep models generally raised accuracy over classical baselines, but at a cost in interpretability: it becomes harder to say why a conversation was flagged \cite{karthikeyan2023}.
%Escalante et al. \cite{escalante2015} argued that a classification based on the whole conversation can only be used for forensics, not prevention, and proposed classifying with partial information instead. 
Vogt et al. \cite{vogt2021} formalised early sexual predator detection, in which a running chat is analysed message by message and a warning should be raised as early and as accurately as possible. Their BERT-based baselines also reached state-of-the-art results for the conventional task. Eilifsen et al. \cite{eilifsen2023} combined a dynamic trust model with a sliding window for early detection, and An et al. \cite{an2025} optimised detection at the turn level using risk labels derived from luring communication theory.






%LLMs have recently been applied to the task. Nguyen et al. (2023) fine-tuned the open-weight Llama 2 7B model to detect predatory chats and abusive language. Hamm and McKeever (2025) fine-tuned the much smaller Llama 3.2 1B model on PAN12 and reported an F1 score of 0.99 for predator identification, outperforming SVMs and earlier models. Kim et al. (2025) used BERT-based models and LLMs for early detection on a Korean version of PAN12, retrieving earlier messages similar to the current one as context instead of using a fixed window.
%Prompted, general-purpose LLMs behave less reliably. Prosser and Edwards (2024a) tested several LLMs across more than 6,000 interactions and found that none was clearly appropriate for grooming prevention. Behaviour was inconsistent, and some models, especially open-source ones, produced potentially harmful answers. In a follow-up, Prosser and Edwards (2024b) found that models detected grooming even in harmless conversations between real people. For an investigative tool this over-flagging is a concrete risk: a ranker that rates innocent segments as suspicious wastes investigator time and can mislead.
%Two conclusions can be drawn. First, conversation- and author-level detection on PAN12 is close to saturated: when fine-tuned models reach F1 scores near 0.99, further gains on that task say little. Second, almost none of this work tells an investigator where in a conversation the evidence lies. That question is the subject of the next section.

\section{Localising evidence within conversations}

\section{Evidence ranking, and LLM as rankers}

\section{LLM faithfulness and hallucination}

\section{Robustness and generalisation}

\section{AI in forensic investigation}

\section{Research gap}